\documentclass[14pt,a4paper]{extarticle}

\usepackage[T2A]{fontenc}
\usepackage[utf8]{inputenc}
\usepackage[russian]{babel}
\usepackage{amsmath,amssymb,amsfonts}
\usepackage{geometry}
\usepackage{graphicx}
\usepackage{float}
\usepackage{booktabs}
\usepackage{array}
\usepackage{hyperref}
\usepackage{xcolor}
\usepackage{caption}

\geometry{
    left=30mm,
    right=15mm,
    top=20mm,
    bottom=20mm
}

\hypersetup{
    colorlinks=true,
    linkcolor=black,
    urlcolor=blue
}

\newcommand{\dd}{\,\mathrm{d}}
\newcommand{\dx}{\Delta x}
\newcommand{\dy}{\Delta y}
\newcommand{\dt}{\Delta t}
\newcommand{\norm}[1]{\left\lVert #1 \right\rVert}

\newcommand{\placeholderfigure}[2]{
\begin{figure}[H]
    \centering
    \fbox{
        \begin{minipage}[c][0.28\textheight][c]{0.86\textwidth}
            \centering
            \vspace{4mm}
            \textbf{Заглушка рисунка}\\[2mm]
            #2\\[4mm]
            Файл будет подставлен после финального численного запуска.
            \vspace{4mm}
        \end{minipage}
    }
    \caption{#1}
\end{figure}
}

\begin{document}

\begin{titlepage}
    \begin{center}
        \small
        Московский государственный университет имени М.~В.~Ломоносова\\
        Механико-математический факультет\\
        Кафедра вычислительной математики

        \vfill

        \Large
        \textbf{Курсовая работа}\\[8mm]

        \large
        \textbf{Численное моделирование двумерного несжимаемого течения\\
        на смещенной сетке и подавление нестационарных гармоник}

        \vfill

        \begin{flushright}
            \normalsize
            Выполнил: студент Илья Макаров\\[2mm]
            Научный руководитель: \underline{\hspace{55mm}}
        \end{flushright}

        \vfill

        Москва, 2026
    \end{center}
\end{titlepage}

\tableofcontents
\newpage

\section*{Введение}
\addcontentsline{toc}{section}{Введение}

В работе рассматривается численное решение двумерных уравнений Навье--Стокса
для несжимаемой вязкой жидкости в прямоугольной области. Расчет строится на
равномерной смещенной сетке: давление хранится в центрах ячеек, а компоненты
скорости расположены на гранях. Такое расположение удобно тем, что разностные
приближения градиента давления и дивергенции скорости строятся на соседних
согласованных контрольных объемах. Это дает более точную разностную
аппроксимацию операторов дифференцирования и позволяет обходиться без лишних
интерполяций между несовпадающими узлами.

Цель работы состоит в том, чтобы:
\begin{enumerate}
    \item реализовать и проверить базовый численный метод на задаче о кювете;
    \item решить задачу с гармонической пульсационной правой частью;
    \item найти нестационарные гармоники около стационарного состояния и
    вычесть соответствующую компоненту из правой части, чтобы получить более
    стационарное течение.
\end{enumerate}

Вычислительный метод использует единую разреженную систему для скоростей и
давления на каждом шаге по времени. Для стационарного режима применяется метод
Ньютона, а для анализа нестационарных гармоник используется спектральная
задача для линеаризованного оператора.

\section{Постановка задачи}

В прямоугольной области
\[
    \Omega = [0,L_x]\times[0,L_y]
\]
решается система двумерных несжимаемых уравнений Навье--Стокса:
\begin{align}
    \frac{\partial \mathbf{u}}{\partial t}
    + (\mathbf{u}\cdot\nabla)\mathbf{u}
    - \nu \Delta \mathbf{u}
    + \nabla p &= \mathbf{f}, \label{eq:ns_momentum}\\
    \nabla\cdot \mathbf{u} &= 0. \label{eq:ns_div}
\end{align}
Здесь
\[
    \mathbf{u}=(u,v), \qquad
    \mathbf{f}=(f_u,f_v),
\]
\(\nu>0\) -- кинематическая вязкость, а \(p\) -- давление, определенное с
точностью до константы. Конкретные значения параметров и правой части
выбираются отдельно для каждой из рассматриваемых задач.

\section{Численная схема и смещенная сетка}

Область разбивается на \(N_x\times N_y\) прямоугольных ячеек с шагами
\[
    \dx=\frac{L_x}{N_x}, \qquad \dy=\frac{L_y}{N_y}.
\]
Давление хранится в центрах ячеек,
\[
    p_{i,j}\approx p\left(\left(i+\frac12\right)\dx,
    \left(j+\frac12\right)\dy\right),
    \quad i=0,\dots,N_x-1,\quad j=0,\dots,N_y-1,
\]
горизонтальная скорость -- на вертикальных гранях,
\[
    u_{i,j}\approx u\left(i\dx,\left(j+\frac12\right)\dy\right),
    \quad i=0,\dots,N_x,\quad j=0,\dots,N_y-1,
\]
а вертикальная скорость -- на горизонтальных гранях,
\[
    v_{i,j}\approx v\left(\left(i+\frac12\right)\dx,j\dy\right),
    \quad i=0,\dots,N_x-1,\quad j=0,\dots,N_y.
\]
Правая часть хранится на тех же узлах, что и соответствующая компонента
скорости: \(f_u\) задается на вертикальных гранях, \(f_v\) -- на
горизонтальных. Из-за такого хранения производная давления по \(x\) сразу
считается в узлах \(u\), производная давления по \(y\) -- в узлах \(v\), а
дивергенция скорости -- в центрах ячеек как баланс потоков через стороны
ячейки.

Вектор неизвестных имеет вид
\[
    X=(u_{\mathrm{int}},v_{\mathrm{int}},p),
\]
где в него входят только внутренние значения скоростей и все значения
давления. Значения скорости на стенках задаются условиями конкретной задачи и
не становятся неизвестными. Если разностный шаблон для вязкого или
конвективного члена касается стенки, известное граничное значение переносится
в правую часть. Давление фиксируется дополнительным условием калибровки
\[
    p_{0,0}=0.
\]

При переходе на следующий временной слой конвективные члены вычисляются по
известному полю скорости, а вязкость, давление и условие несжимаемости
берутся на новом слое. В матричной форме один шаг записывается как
\[
    \begin{pmatrix}
        A_u & 0 & G_x\\
        0 & A_v & G_y\\
        D_x & D_y & 0
    \end{pmatrix}
    \begin{pmatrix}
        u^{n+1}\\
        v^{n+1}\\
        p^{n+1}
    \end{pmatrix}
    =
    \begin{pmatrix}
        r_u(u^n,v^n,\mathbf{f})\\
        r_v(u^n,v^n,\mathbf{f})\\
        0
    \end{pmatrix}.
    \label{eq:monolithic_step}
\]
Здесь \(G_x,G_y\) -- дискретный градиент давления, \(D_x,D_y\) -- дискретная
дивергенция, а \(A_u,A_v\) содержат вклад временной производной и вязкого
члена. Матрица зависит только от сетки, шага по времени, вязкости и способа
задания стенок, поэтому для фиксированного расчета она собирается один раз.
После этого на каждом шаге меняется только правая часть: в нее входят
предыдущее поле скорости, конвекция, внешняя сила и известные добавки от
стенок. Полученная разреженная система решается прямым методом для разреженных
матриц.

Для поиска стационарного решения используется отображение \(\Phi\), которое
одному полю скорости ставит в соответствие результат одного шага
\eqref{eq:monolithic_step}. Стационарное поле находится из уравнения
\[
    \Phi(U)-U=0
\]
методом Ньютона. На каждой итерации линейная система для поправки решается
итерационным методом GMRES, а произведение якобиана на вектор вычисляется
конечной разностью через повторный вызов шага по времени. Для устойчивости
обновления применяется демпфирование шага.

После нахождения стационарного состояния строится линеаризация около него.
Полученный оператор действует только на скоростные возмущения, а результат
проецируется на бездивергентные поля через ту же монолитную систему для
скорости и давления. Собственные значения и собственные векторы ищутся
методом Арнольди без явной сборки полной плотной матрицы. Для дальнейшего
анализа выбираются моды с наибольшей вещественной частью собственного
значения: именно они сильнее всего отвечают за нестационарное поведение около
найденного стационара.

\section{Задача о кювете}

Задача о кювете используется как базовая проверка численного метода. В этой
постановке правая часть равна нулю, а движение создается верхней крышкой
кюветы. На верхней стенке задается
\[
    u=1,\qquad v=0,
\]
на всех остальных стенках обе компоненты скорости равны нулю. В результате
ожидается стационарная циркуляция внутри области, по которой можно проверить
качественную корректность поля скорости, давления и вихря.

\placeholderfigure{Поле скорости в задаче о кювете.}{
    Ожидаемая замена: \texttt{plots/run/final\_state/streamlines.png}
}

\placeholderfigure{Давление и вихрь в задаче о кювете.}{
    Ожидаемая замена: \texttt{plots/run/final\_state/pressure.png} и
    \texttt{plots/run/final\_state/vorticity.png}
}

\section{Задача с пульсационной правой частью}

Во второй задаче стенки неподвижны:
\[
    u=0,\qquad v=0
    \quad \text{на всей границе области}.
\]
Движение создается ненулевой правой частью. В текущем конфигурационном файле
задана только горизонтальная компонента силы, а вертикальная компонента равна
нулю:
\[
    f_v(x,y)=0,
\]
\begin{align}
    f_u(x,y)={}&
    -\frac{19}{4\pi}\sin(2\pi x)\cos(4\pi y)
    -\frac{3}{2\pi}\sin(4\pi x)\cos(4\pi y) \notag\\
    &+\frac{7}{4\pi}\sin(6\pi x)\cos(4\pi y)
    +\frac{7}{4\pi}\sin(2\pi x)\cos(8\pi y) \notag\\
    &-\frac{19}{200\pi}\sin(\pi x)\cos(6\pi y).
    \label{eq:forcing}
\end{align}
Эта сила является суммой нескольких пространственных гармоник. Нестационарный
расчет показывает, какие структуры возникают под действием такого внешнего
возбуждения и насколько быстро решение приближается к стационарному режиму.

\placeholderfigure{Поле скорости в задаче с пульсационной правой частью.}{
    Ожидаемая замена: \texttt{plots/run/final\_state/streamlines.png}
}

\placeholderfigure{Давление и вихрь в задаче с пульсационной правой частью.}{
    Ожидаемая замена: \texttt{plots/run/final\_state/pressure.png} и
    \texttt{plots/run/final\_state/vorticity.png}
}

\section{Стационарное состояние и подавление пульсаций}

После расчета с пульсационной правой частью находится стационарное состояние,
соответствующее той же постановке. Затем около этого состояния строится
линейная модель для малых возмущений и находятся собственные моды. Моды с
наибольшей вещественной частью собственного значения выделяют направления,
которые поддерживают нестационарное поведение.

Чтобы ослабить эти пульсации, правая часть раскладывается на компоненту вдоль
выбранных собственных мод и остаток. В дальнейший расчет передается
отфильтрованная правая часть
\[
    F_{\mathrm{filtered}} = F - F_{\mathrm{osc}},
\]
где \(F_{\mathrm{osc}}\) -- проекция исходного возбуждения на найденное
пульсационное подпространство. После этого выполняется повторный запуск и
сравнивается поведение решения до и после вычитания.

\placeholderfigure{Стационарное поле скорости.}{
    Ожидаемая замена: \texttt{plots/steady/streamlines.png}
}

\placeholderfigure{Спектр линеаризованного оператора около стационара.}{
    Ожидаемая замена: график собственных значений из
    \texttt{plots/linearized/eigenpairs.pkl}
}

\placeholderfigure{Сравнение расчета до и после вычитания пульсационных мод.}{
    Ожидаемая замена: \texttt{plots/projected\_run/stokes\_velocity\_change.png}
}

\placeholderfigure{Финальное поле после вычитания пульсационных мод.}{
    Ожидаемая замена: \texttt{plots/projected\_run/final\_state/streamlines.png},
    \texttt{pressure.png}, \texttt{vorticity.png}
}

\section*{Заключение}
\addcontentsline{toc}{section}{Заключение}

В работе описан численный метод для двумерных несжимаемых уравнений
Навье--Стокса на смещенной сетке. Давление хранится в центрах ячеек, скорости
-- на гранях, а правая часть задается в тех же узлах, что и соответствующие
компоненты скорости. Такое хранение дает более точные и согласованные
разностные аппроксимации градиента и дивергенции и позволяет на каждом шаге
решать одну монолитную систему для скорости и давления.

Метод проверяется на задаче о кювете, затем применяется к задаче с
гармонической правой частью. Для этой задачи находится стационарное состояние,
строится спектральная задача для малых возмущений и выделяются моды, связанные
с нестационарным поведением. Вычитание соответствующей компоненты из правой
части позволяет получить более стационарное течение при той же основной
постановке.

\end{document}
